\textbf{РЕЦЕНЗИЯ}\textbf{}

На выпускную квалификационную работу бакалавра

Вылегжаниной Алёны Игоревны

"Метод восстановления глубины архивных фотографий по закраске изображения"

В современном мире активно развиваются технологии трехмерной реконструкции и визуализации, что открывает новые возможности для сохранения и воссоздания культурного наследия. Особую ценность представляет возможность преобразования архивных двумерных фотографий в объемные трехмерные модели, что позволяет более полно восстановить историческую информацию, улучшить восприятие исторических объектов и создать их виртуальные реконструкции.

В данной работе был разработан метод восстановления трехмерной глубины из архивных фотографий на основе гибридного подхода, объединяющего нейросетевой метод и классический алгоритм восстановления структуры по закраске.

В аналитической части проведен обзор существующих методов преобразования двумерных изображений в трехмерные, выполнена классификация методов по нескольким критериям. Проанализированы различные подходы к восстановлению трехмерной информации.

В конструкторской части представлена структура разработанного метода восстановления трехмерной информации из архивных фотографий.

В технологической части реализован программный комплекс на языке Python с использованием специализированных библиотек и проведено его тестирование на различных типах изображений.

В исследовательской части выполнено сравнение разработанного метода с существующими подходами, определены его преимущества и ограничения.

К достоинствам разработанного метода можно отнести то, что в большинстве случаев модифицированный метод показал более точный результат, чем базовые по ключевым метрикам.

К недостаткам разработанного метода можно отнести то, что в некоторых случаях базовые методы показали более точные результаты, чем модифицированный по ключевым метрикам.

Несмотря на отмеченные недостатки, выпускная квалификационная работа Вылегжаниной А.И. соответствует квалификационным требованиям, предъявляемым к работе бакалавра, и заслуживает отличной оценки, а автор – присвоения степени бакалавра по направлению 09.03.04 «Программная инженерия».

Рецензент:

Д. т. н., доцент ИУ6 _______________________________\underline{Попов А. Ю.}

29.05.2025 г.

